% Part: propositional-logic
% Chapter: syntax-and-semantics
% Section: semantic-notions

\documentclass[../../../include/open-logic-section]{subfiles}

\begin{document}

\olfileid[pt]{pl}{syn}{sem}

\olsection{Noções semânticas}

Definimos as seguintes noções semânticas:

\begin{defn} 
\begin{enumerate}
\item !!^a{formula}~$!A$ é \emph{satisfatível} se, para
  algum~$\pAssign{v}$, $\pSat{v}{!A}$; ela é
  \emph{insatisfatível} se não há $\pAssign{v}$ tal que $\pSat{v}{!A}$;
\item !!^a{formula}~$!A$ é uma \emph{tautologia} se $\pSat{v}{!A}$ para
  todas as !!{valuation}s~$\pAssign{v}$;
\item !!^a{formula}~$!A$ é \emph{contingente} se é satisfatível, mas
  não é uma tautologia;
\item Se $\Gamma$ é um conjunto de !!{formula}s, $\Gamma \Entails !A$ (``$\Gamma$
  tem $!A$ como consequência'') se, e somente se, $\pSat{v}{!A}$ para toda
  !!{valuation}~$\pAssign{v}$ para a qual $\pSat{v}{\Gamma}$.
\item Se $\Gamma$ é um conjunto de !!{formula}s, $\Gamma$ é
  \emph{satisfatível} se existe !!a{valuation}~$\pAssign{v}$ para a qual
  $\pSat{v}{\Gamma}$, e $\Gamma$ é
  \emph{insatisfatível} caso contrário.
\end{enumerate} 
\end{defn}

\begin{prob} 
 Para cada uma das quatro !!{formula}s a seguir, determine se ela é
 (a)~satisfatível, (b)~uma tautologia e (c)~contingente.
\begin{enumerate}
  \item \( ( \Obj p_0 \lif ( \lnot \Obj p_1 \lif \lnot \Obj p_0 ) ) \).
  \item \( ( ( \Obj p_0 \land \lnot \Obj p_1 ) \lif ( \lnot \Obj p_0 \land \Obj p_2 )) \liff ( ( \Obj p_2 \lif \Obj p_0 ) \lif ( \Obj p_0 \lif \Obj p_1 )) \).
  \item \( ( \Obj p_0 \liff \Obj p_1 ) \lif ( \Obj p_2 \liff \lnot \Obj p_1 ) \).
  \item \( (( \Obj p_0 \liff ( \lnot \Obj p_1 \land \Obj p_2 )) \lor ( \Obj p_2 \lif ( \Obj p_0 \liff \Obj p_1 ))) \).
\end{enumerate}
\end{prob}

\begin{prop}
\ollabel{prop:semanticalfacts} 
\begin{enumerate} 
\item $!A$ é uma tautologia se, e somente se,
  $\emptyset \Entails !A$; 
\item Se $\Gamma \Entails !A$ e $\Gamma \Entails !A \lif !B$, então
  $\Gamma \Entails !B$;
\item Se $\Gamma$ é satisfatível, então todo subconjunto finito de $\Gamma$
  também é satisfatível; 
\item \ollabel{def:monotonicity} Monotonicidade: se $\Gamma \subseteq \Delta$
  e $\Gamma \Entails !A$, então também $\Delta \Entails !A$;
\item \ollabel{def:Cut} Transitividade: se $\Gamma \Entails !A$ e
  $\Delta \cup \{ !A\} \Entails !B$, então $\Gamma \cup \Delta \Entails
  !B$.
\end{enumerate}
\end{prop}

\begin{proof}
Exercício.
\end{proof}

\begin{prob}
Demonstre \olref[pl][syn][sem]{prop:semanticalfacts}
\end{prob}

\begin{prop}\ollabel{prop:entails-unsat}
  $\Gamma \Entails !A$ se, e somente se, $\Gamma \cup \{\lnot !A\}$
  é insatisfatível.
\end{prop}

\begin{proof}
Exercício.
\end{proof}

\begin{prob}
Demonstre \olref[pl][syn][sem]{prop:entails-unsat}
\end{prob}

\begin{thm}[Teorema semântico da dedução]
  \ollabel{thm:sem-deduction} $\Gamma \Entails !A \lif !B$ se, e somente
  se, $\Gamma \cup \{!A\} \Entails !B$.
\end{thm}

\begin{proof}
Exercício.
\end{proof}

\begin{prob}
Demonstre \olref[pl][syn][sem]{thm:sem-deduction}
\end{prob}

\end{document}
